\documentclass[class=scrbook, crop=false]{standalone}
\usepackage[subpreambles=true]{standalone}
\ifstandalone
    \input{../settings+/settings}
\fi

% ----------------------------------------------------------------------------
%                                Introduction
% ----------------------------------------------------------------------------
\begin{document}

\chapter{Introduction}
\label{Chapter::Introduction}

\section{Motivation}
\label{sec:motivation}

In the field of optimal control, 
direct multiple shooting~\cite{bock1984multiple} is a widely used method for solving continuous-time optimal control problems,
which transcribes the OCP into a nonlinear programming problem by discretizing the time horizon into multiple shooting intervals~\cite{Diehl2018}.

This NLP is commonly solved by SQP, which computes each step from a QP subproblem 
which inherits an optimal control structure that Riccati recursion can exploit \cite{rao1998application}.
To retain the fast local convergence of SQP, 
the QP subproblems are often formulated with the exact Hessian of the Lagrangian \cite{BoggsTolle1995},
this could result in an indefinite Hessian thus a nonconvex QP as subproblem, 
which is not guaranteed to have a unique solution and can cause the Riccati recursion to fail.

A common remedy is to regularize at the SQP level,
i.e.\ to modify the QP Hessian before the QP is passed to the solver.
Such a correction is decided from the Hessian of the QP subproblem alone and therefore overlooks two aspects.
First, inside an interior point method, the barrier term of the inequality constraints adds curvature to the Hessian,
which can render the system well-defined even when the QP Hessian itself is indefinite.
Second, when the correction targets the full Hessian, it ignores that positive definiteness of the reduced Hessian
is already sufficient for the QP to have a unique solution.

This thesis therefore regularizes at the QP level, inside the solver,
where the correction is computed from the factorized matrices and applied only when indefiniteness is actually detected.

\section{Contributions}
\label{sec:contributions}
The main contributions of this thesis are the following:
\begin{enumerate}
    \item \textbf{Riccati-based primal regularization}, which regularizes the reduced Hessian blocks
    obtained from the Riccati recursion, only at the stages where the Cholesky factorization fails
    (Section~\ref{sec:riccati_primal_reg}).

    \item \textbf{An improved stage regularization strategy}, derived from showing that minimum perturbation
    propagates negative curvature through the cost-to-go matrix to the preceding stages
    (Section~\ref{sec:fail_minimal_perturbation}).

    \item \textbf{A proof of asymptotic inactivity}, so that the local convergence rate of the interior point method is preserved
    (Section~\ref{sec:asymptotic_inactivity}).

    \item \textbf{A benchmark of $175$ indefinite OCP QPs} collected from the SQP iterations of eight OCPs,
    with an analysis of where the indefiniteness originates (Chapter~\ref{Chapter::Benchmark}).

    \item \textbf{A numerical comparison} on this benchmark, in which the best Riccati-based method leaves $4$ instances unsolved,
    against $32$ to $59$ for full Hessian regularization and $15$ for inertia correction
    (Section~\ref{sec:comparison}).
\end{enumerate}

\section{Related Solvers}
\label{sec:related_solvers}
This thesis builds on two open-source software packages for embedded optimal control, \texttt{acados} and HPIPM.

\texttt{acados}~\cite{verschueren2020acadosmodularopensourceframework} is a modular framework for fast NMPC.
It transcribes the OCP by direct multiple shooting, solves the resulting NLP by SQP or its real-time iteration variant,
and passes the QP subproblems to one of several interfaced QP solvers.
When the exact Hessian is used, \texttt{acados} handles indefiniteness at the SQP level:
the QP Hessian can be regularized before the QP is passed to the solver or be replaced by a Gauss--Newton approximation.
In this thesis, \texttt{acados} is used to generate the SQP iterates from which the benchmark QPs are collected (Section~\ref{sec:benchmark_construction}).

HPIPM~\cite{frison2020hpipmhighperformancequadraticprogramming} is a primal--dual interior point solver
for OCP-structured QPs and the default QP solver in \texttt{acados}.
It solves the KKT system of each iteration by Riccati recursion, with a cost linear in the horizon length,
and relies on high-performance linear algebra routines tailored to the small and medium-sized matrices of optimal control.
HPIPM assumes a convex QP: the Riccati recursion requires the reduced Hessian blocks to be positive definite, so it fails on indefinite QPs.
This failure is precisely the criterion by which the benchmark QPs are collected (Section~\ref{sec:benchmark_construction}),
and the prototype solver used in the experiments reproduces the iteration behavior of HPIPM on convex QPs, 
so that the regularization is the only difference between the two (Section~\ref{sec:prototype}).


\section{Outline}
\label{sec:outline}
The remainder of this thesis is organized as follows.
Chapter~\ref{Chapter::Theoretical_Background} introduces the primal--dual interior point method for OCP-structured QPs,
the Riccati recursion used to solve its KKT system,
and the barrier update and step size control strategies adopted in the solver.
Chapter~\ref{Chapter::Primal_Regularization} classifies primal regularization by its target matrix,
develops the Riccati-based primal regularization,
derives the improved stage regularization strategy from the failure of minimum perturbation,
and proves that the regularization is asymptotically inactive.
Chapter~\ref{Chapter::Benchmark} constructs the benchmark of indefinite QPs from the SQP iterations of eight OCPs
and analyzes how the indefiniteness is distributed over the stages and where it originates.
Chapter~\ref{Chapter::Experiments} describes the prototype solver and the experiment setup,
compares the Riccati-based methods with full Hessian regularization and inertia correction on the benchmark,
and discusses the instances on which the methods fail.
Finally, Chapter~\ref{Chapter::Conclusion} summarizes the results and outlines directions for future work.

% You can use this to add content for standalone documents if you like
% In this case, we would like to show the acronyms and references.
\ifstandalone
    \printnoidxglossary[type=\acronymtype]
    \printbibliography[heading=bibintoc]
\fi

\end{document}
